\chapter{Introduction}
\label{ch:introduction}

This chapter provides an overview of the thesis...

% This gives the reader useful information to understand the problem
Verification of neural networks is very important...

% This tells the reader the problem your solution addresses
However, larger networks can be hard to process...

% This tells the reader what you did
We show a simple method to simplify big networks...

% This tells the reader why what you did is correct
We show that our simplification is sound...

% This tells the reader the organisation of the thesis, namely what sections are there and what they contain
\section{Outline}
\label{sec:outline}
This section outlines the organization of the thesis:
\begin{itemize}
	\item \textbf{\Cref{ch:background}}: Introduces key concepts to understand the contributions...
	\begin{itemize}
		\item \textbf{\Cref{sec:neural-networks}}: Introduces the concepts of neural networks and deep learning...
		\item \textbf{\Cref{sec:interaction-nets}}: Introduces the interaction net computational model...
		\item \textbf{\Cref{sec:satisfiability-modulo-theories}}: Introduces SMT solvers...
	\end{itemize}
	\item \textbf{\Cref{ch:core}}: Dives deeper into the implementation...
	\item \textbf{\Cref{ch:related-work}}: Analyzes existing verification approaches...
	\item \textbf{\Cref{ch:conclusion}}: Summarizes the findings and discusses future work...
\end{itemize}
